\section{Exact contracts and their exceptions}\label{s5:contracts}

Exact external contracts can remove aggregate equations that obstruct
path elimination. They lead to several different algorithms, whose
hypotheses and certificate guarantees should be distinguished. Throughout
this section there is one pool, no pool--pool arc, and finite rational
bounds on every actual flow. Write $C_i=\lambda_i$ for an input quality,
$y_i,v_j,z_{ij}$ for feed, outlet, and bypass flows, and $q$ for pool
quality. An absent feed or outlet has bounds $[0,0]$. Exact supplies are
$a_i$, and an exact product contract is demand $b_j$ and quality $B_j$.
The symbols $C_i,B_j,q$ denote vectors except in explicitly scalar
arguments. All lower bounds, including positive individual arc bounds,
are retained. Exact quality at zero demand has the homogeneous meaning
of Section~\ref{s1:model}.

The main reductions use contracts in three ways. Fixed outlet sets
concentrate the conserved input mass and attribute mass into a bounded
boundary. Full contracts turn scalar quality equations into signed
node divergences, whose feasibility cuts can be enumerated on paths
and cycles. With a fixed number of exceptions, an active ordinary
outlet confines the entire pool-quality vector to an affine plane,
even if the ambient quality rank grows. Restrictive common capacity
and two distinct full input vectors require separate arguments below;
neither conclusion follows from the conservation identities alone.

The algorithms use rational LP, exact arithmetic in a common represented
real-algebraic field, and fixed-dimensional real-algebraic decision and
sampling. The latter includes polynomially growing degrees under dense
encoding: see \citet[Theorem~1.3.1 and Sections~3.1.3--3.2]
{basu1996-on-the-combinatorial-and-algebraic}. The algebraic LP step uses
\citet[Section~5, Remark~1]{s4:adler1994}. Polynomial time below always
means rational bit time, with an exponent allowed to depend on the
specified fixed parameters.

\subsection{Parameterized path relations}\label{s5:parameter-paths}

We first retain two general projection facts. They concern the relation
on retained coordinates; they do not eliminate arbitrary accumulated
costs or dense aggregate constraints.

\begin{theorem}[One parameter and general planar relations]
\label{s5:quasipoly-path}
Let $\rho$ range over a compact rational interval and let
$(x_{i-1},x_i)\in P_i(\rho)$, $1\le i\le h$, where each $P_i$ is
specified by weak inequalities linear in its two state coordinates,
with explicitly encoded rational polynomial coefficients of fixed
degree in $\rho$. Suppose every state has a uniform finite rational
box bound. For input length $N$, the exact endpoint relation has a
representation on $N^{O(\log(h+1))}$ parameter cells. Each cell is an
open interval or singleton with polynomial-degree, polynomial-encoding
algebraic endpoints. On each cell the relation has polynomially many
linear endpoint inequalities with polynomial-degree, polynomial-bit
integer-polynomial coefficients. This representation is constructible
in quasipolynomial bit time. Existential feasibility has the same bound
and admits a polynomial-encoding witness in one polynomial-degree real
algebraic field.
\end{theorem}
\begin{proof}
Clear rational denominators in the initial rows. Include both endpoints
of the supplied parameter interval as singleton cells. Compose along a
balanced binary tree, padding with bounded identity relations if necessary.
Consider two children on a common open parameter cell, each with a fixed
row list whose coefficients have degree at most $D$. In coordinates
$(x,y,z)$, eliminate $y$ by Fourier--Motzkin elimination. For rows with
$y$ coefficients $b_i>0>b_j$, use $-b_j$ times the first row plus $b_i$
times the second. Retain all zero-$y$ rows and all applicable pairs,
including pairs from the same child. There are quadratically many
candidates in the total child row count, and coefficient degrees at most
$2D$. Split first at roots of the $b_i$ to fix the applicable pairs.
At a root, treat a vanishing coefficient as zero, using its actual value.

For the resulting rows $A_rx+B_rz\le C_r$, retain the four fixed box
rows and split further at every root of every nonzero minor of orders
one, two, and three in the augmented matrix $(A_r,B_r,C_r)$. For $k$
candidates there are $O(k^3)$ such polynomials, of degree $O(D)$.
Their signs fix whether any two boundary lines intersect uniquely and,
by Cramer's rule, whether that intersection satisfies each row. Every
nonempty bounded planar polyhedron, including a point or segment, has
a vertex determined by two independent boundary lines.

On a sample in each open sign cell, delete redundant candidate rows
one at a time, keeping the box. The same subset remains equivalent throughout the
cell: each potential vertex of its bounded polyhedron has invariant
feasibility and invariant satisfaction of every omitted row. Thus
validity at all vertices persists. Emptiness also persists, since the
existence of a feasible intersection pair is sign-determined. At every
singleton perform the same selection in the field of that parameter.
In particular, recompute zero coefficients and applicable elimination
pairs there; specializing either neighboring open-cell description
need not give the singleton projection. Store selected symbolic rows,
rather than substituting unrelated algebraic coefficients into subsequent
steps. Each row is required to represent the child only on its stored
cell, which suffices because a parent uses it only on a common refinement.
On an empty sign cell
store the constant contradiction $0\le-1$ together with the box,
so an empty relation also has a constant-size retained description.

Lemma~\ref{s4:polygon-composition} bounds the necessary number of rows
in a nonempty composition by a universal constant times the sum of
child sizes. This also bounds a subset of the candidate description.
In dimension two keep the facet rows. For a segment, rows active at a
relative interior point have normals perpendicular to its line; there
must be normals of both signs, since otherwise a small displacement
off the line would remain feasible. Two such rows define the affine
hull, and at most two more rows delimit its endpoints. For a point,
the active normals positively span the plane, and at most four suffice:
either three surround the origin or two opposite pairs do. If the
active normals did not positively span, a nonzero feasible displacement
would contradict that the set is a point. Thus in each dimension a
bounded-size subset of the original candidate rows gives the required
description. Retaining the four box rows only adds four to these bounds.
Thus, with a larger universal $C$ and the retained box,
\[
 m_{\rm parent}\le C(m_{\rm left}+m_{\rm right}+2).
\]
Logarithmic depth makes all row counts polynomial in $N$.

In one parameter the common refinement of the two child partitions has
only the sum of their numbers of breakpoints. On each common cell,
coefficient roots and minor roots introduce polynomially many additional
cells. Consequently
\[
 L_{\rm parent}\le N^c(L_{\rm left}+L_{\rm right})
\]
for a fixed $c$. Iteration through $O(\log(h+1))$ levels proves the
quasipolynomial bound, including all intermediate work. It is essential
here that the parameter is scalar.

Each coefficient combination takes two polynomial products and a sum.
Degree at most doubles per level, so it is $O(dh)$ for initial degree
$d$. Coefficient bit lengths satisfy the same doubled recurrence with an
additional logarithmic convolution term and remain polynomial in $N$.
All minors therefore have polynomial degree and coefficient bit length. Exact univariate
root isolation, comparison, and sign evaluation have polynomial cost
per polynomial; separation bounds give polynomial-bit rational samples
between distinct roots. Singleton arithmetic takes place in the field
of the single chosen parameter and introduces no new parameter root.

A final minor refinement makes endpoint nonemptiness constant even for
$h=1$. Check every open-cell sample and singleton. At a feasible
parameter, solve the original bounded path LP over its field. A vertex
basis expresses its coordinates as Cramer ratios of the original
polynomial coefficients evaluated at that parameter. The determinants
have polynomial degree and coefficient bit length: a determinant of dimension $O(N)$
contains at most $O(N)!$ terms, so its coefficient bit length increases by
only polynomially many bits. These expressions give a polynomial-size
common-field witness; the algebraic LP algorithm constructs one.
Alternatively, store the balanced composition trees and lift interval
slices as in Theorem~\ref{s4:path-projection}.
\end{proof}

\begin{theorem}[Affine strips, trees, and fixed total cycle rank]
\label{s5:strip}
Let $\eta$ have fixed dimension and range over a supplied
quantifier-free rational polynomial domain, possibly all of its ambient
space. Suppose a scalar path has finite polynomial coordinate bounds
at every parameter and transitions
\[
 \ell_i(\eta)\le x_i\le u_i(\eta),\qquad
 a_i(\eta)x_{i-1}+b_i(\eta)\le x_i
 \le a_i(\eta)x_{i-1}+c_i(\eta).
\]
All polynomials are densely encoded; degrees may grow. On each realizable
sign condition of the gains $a_i$, the projection onto $\eta$ and a fixed
number of states has polynomial size, degree, and coefficient length.
The same holds on a tree with arbitrarily oriented edges, and on a graph
of fixed total cycle rank after retaining at most twice that many extra
states. These constructions give polynomial-time feasibility and exact
infimum and attainment analysis for polynomial objectives on the
retained coordinates. An optimizer is returned when the infimum is
attained; a nonempty compact closed original feasible set suffices.
\end{theorem}
\begin{proof}
First consider a path and a sign cell where every gain is nonzero. Put
$A_0=1$, $A_i=\prod_{h=1}^i a_h$, and $y_i=x_i/A_i$. The transitions
become $\alpha_i\le y_i-y_{i-1}\le\beta_i$, where the endpoints
$b_i/A_i,c_i/A_i$ are ordered according to the known sign of $A_i$.
The state bounds similarly become $L_i\le y_i\le U_i$. Include every
lower-versus-upper consistency condition. Define directed distances
\[
 d(i,j)=\sum_{h=i+1}^j\beta_h,\qquad
 d(j,i)=-\sum_{h=i+1}^j\alpha_h\quad(i<j),\qquad d(i,i)=0.
\]
These are shortest-path distances: every walk on the bidirected path
reduces to its unique simple path and two-edge excursions of lengths
$\beta_h-\alpha_h\ge0$. Feasibility is exactly
\begin{equation}
 L_i\le U_j+d(j,i)\quad\hbox{for every }i,j.
 \label{s5:distance-test}
\end{equation}
Necessity follows by summing differences. For sufficiency take
$y_i=\min_j(U_j+d(j,i))$. This lies between $L_i$ and $U_i$; the
triangle inequality yields $y_i\le y_h+d(h,i)$ and hence both
adjacent transition inequalities.

To retain a state, replace its interval by its prospective exact value
$x_i/A_i$ and keep its original bounds separately. This gives
$O(n^2)$ inequalities on the retained coordinates. Multiple interval
bounds are handled by all lower/upper pairs, so no choice of an active
bound is required. A zero gain gives simply $b_i\le x_i\le c_i$ and
no restriction on $x_{i-1}$: split there, place that interval on its
original head, and apply the construction separately.

For a tree, delete zero-gain edges in precisely this way. On each
remaining tree choose factors satisfying $A_v=a_eA_u$ for each original
edge $u\to v$. Rooting assigns factors as products of gains or their
reciprocals along root paths. Dividing by $A_v$ gives the same two
bounded differences on each edge, with signs ordered using $A_v$.
The unique simple path between any two vertices again gives distances
and \eqref{s5:distance-test}, with exactly the same recovery proof.
No zero-gain equation is reversed when an edge points toward the root.

Products of polynomially many densely encoded gains have polynomial
degree and coefficient bit length in a fixed number of parameters. A common product of
all nonzero gains, with a fixed sufficient multiplicity and the original
coefficient denominators, clears all the products, reciprocals, and
path sums just used. Its sign is known on the cell. Clearing denominators
in \eqref{s5:distance-test} therefore gives polynomially encoded rows.
Enumerating realizable gain sign conditions has polynomial complexity
in fixed dimension; zero and lower-dimensional conditions are included.
Recovery uses minima and rational functions in the same field as the
retained sample.

Finally, remove $c$ edges to form a spanning forest when the total cycle
rank is $c$. Retain their at most $2c$ endpoints, project the forest,
and restore every removed original strip row on these states. The
remaining dimension is fixed. Quantifier elimination applied to the
attainable objective values gives the stated feasibility, infimum, and
attainment conclusions. Finite state bounds at each parameter do not
alone imply bounded parameters or attainment; the theorem makes no such
inference.
\end{proof}

Both sides of a strip must have the same gain. Arbitrary planar
relations instead satisfy Theorem~\ref{s5:quasipoly-path}, whose bound
is quasipolynomial and is not asserted to be tight. Neither theorem
retains an arbitrary dense objective. For scalar upper product quality
bounds alone, lowering a fixed pool quality preserves each local
quality row because its pool flow coefficient is nonnegative. Thus
local path feasibility is nested in that parameter. This does not
control changing facets or the global pool mixing equality, which is
not a local upper-bound row. Likewise, inequalities
$x_i\ge|x_{i-1}-a|$ do not enforce iterated absolute values: existential
intermediate states may exceed their lower bound. Such observations
cannot replace either projection proof.

\subsection{A fixed number of pool outlets}\label{s5:fixed-outlets}

\begin{theorem}[Conserved boundary masses]\label{s5:fixed-products}
Suppose the bypass graph has maximum degree two, every source has exact
supply $a_i$, and the set $J_P$ of allowed pool outlets has fixed size
$r$. Every product outside $J_P$ has an exact product contract; products
in $J_P$ may have arbitrary demand and quality intervals. Feasibility
and standard economic optimization are polynomial-time solvable, with
an exact algebraic original optimizer when feasible. Arbitrarily many
feeds and quality coordinates, arbitrary quality rank, restrictive
common pool bounds, and all individual arc bounds are allowed. Linear
costs on the $r$ outlets and on pool throughput are also permitted.
\end{theorem}
\begin{proof}
If $r=0$, mass conservation forces every feed to zero. Check zero against
all incident and common bounds and solve the remaining rational LP,
retaining external contracts. Otherwise eliminate every feed by
\[
 Y_i=a_i-\sum_jz_{ij}.
\]
Put its actual feed bounds on $Y_i$, or impose $Y_i=0$ if absent. Each
row uses at most two bypass coordinates. At $j\notin J_P$ impose
$\sum_i z_{ij}=b_j$ and $\sum_i C_i z_{ij}=b_jB_j$. These also have
at most two coordinates, regardless of the number of attributes.

Retain the bypass set $H$ entering $J_P$, of size at most $2r$. Define
\begin{align}
 T&=\sum_i a_i-\sum_{j\notin J_P}b_j-\sum_{ij\in H}z_{ij},\notag\\
 Q&=\sum_i C_i a_i-\sum_{j\notin J_P}b_jB_j
                         -\sum_{ij\in H}C_i z_{ij}.
 \label{s5:boundary-masses}
\end{align}
Summing the exact source and nonreceiving-product equations gives,
for every local bypass lift,
\begin{equation}
 \sum_iY_i=T,\qquad \sum_i C_iY_i=Q.
 \label{s5:boundary-identities}
\end{equation}
Their coefficients are polynomial-bit sums of rational products.
Retain $v_j$ and outlet fractions $\theta_j$ for all $j\in J_P$, and
require
\[
 \theta_j\ge0,\quad\sum_{j\in J_P}\theta_j=1,\qquad
 v_j=\theta_jT.
\]
There are at most $4r$ retained coordinates. Keep all outlet and common
pool bounds, every receiving-product demand bound on
$t_j=v_j+\sum_i z_{ij}$, and every quality bound on the mass
$\theta_jQ+\sum_i C_i z_{ij}$. These are quadratic rows in the core;
additional quality coordinates add rows, not variables.

Delete $J_P$ from the bypass graph. Every residual component with a
boundary is a path with at most two retained endpoint arcs, and
Theorem~\ref{s4:path-projection} gives its polynomial-size rational
endpoint relation. A component with one boundary gives an interval;
a component with no boundary, including a cycle or isolated node, is
a bounded rational LP. Such a detached component may still feed the
pool. Its arbitrary feasible lift does not change
\eqref{s5:boundary-identities}, so it needs no extra aggregate variable.
Two boundary arcs may meet the same deleted product.

The resulting core is closed, bounded, fixed-dimensional, and defined
by polynomially many quadratic or linear rows of polynomial bit length.
If $T>0$, \eqref{s5:boundary-identities} gives actual quality $Q/T$;
its mass into product $j$ is $v_jQ/T=\theta_jQ$. If $T=0$, nonnegative
$Y_i$ all vanish, so $Q=0$ and every outlet vanishes. Fractions then
have no physical effect. Conversely a feasible physical solution gives
$\theta_j=v_j/T$ when $T>0$ and any simplex point when $T=0$.
Thus core feasibility and original feasibility are equivalent, including
inactive pools and zero-demand products.

For unit input costs $c_i$ and revenues $R_j$, profit is
\[
 \sum_{j\notin J_P}R_jb_j-\sum_ic_i a_i
       +\sum_{j\in J_P}R_j\Bigl(v_j+\sum_i z_{ij}\Bigr).
\]
It is affine in the core. Outlet costs and a cost proportional to $T$
are affine there as well. Fixed-dimensional algebraic optimization gives
an attained optimum, and stored path relations lift its sample in the
same represented field. Recover feeds from $Y_i$ afterward.
\end{proof}

Arbitrary intake costs become a dense bypass objective after the
substitution for $Y_i$; arbitrary bypass costs are already dense. They
are outside Theorem~\ref{s5:fixed-products}. A fixed number of additional
cost-bearing arcs can be retained using Theorem~\ref{s4:fixed-support};
for a designated intake retain its at most two adjacent bypass flows.
Exact nonreceiving quality is material. For example, two unit supplies
have qualities zero and two. Only the clean input bypasses to a product
of demand one and upper quality one; the other product receives one unit
from the pool and has upper quality $3/2$. The first product consumes
all clean supply, forcing pool quality two, hence infeasibility. Treating
its upper quality one as an equality would incorrectly subtract one
unit of quality mass in \eqref{s5:boundary-masses}.

\subsection{Fully contracted signed flows}\label{s5:scaled}

\begin{lemma}[Contract conservation]\label{s5:conservation}
With exact source supplies and exact product contracts, first check
\begin{equation}
 \sum_i a_i=\sum_j b_j,\qquad
 \sum_i C_i a_i=\sum_j B_jb_j.
 \label{s5:global-constants}
\end{equation}
After eliminating $Y_i=a_i-\sum_jz_{ij}$ and
$V_j=b_j-\sum_i z_{ij}$, retain their actual bounds and every equation
\begin{equation}
 \sum_i(C_i-q)z_{ij}=b_j(B_j-q).
 \label{s5:contract-quality}
\end{equation}
Together with bypass bounds, these conditions are equivalent to all
physical constraints except a possibly restrictive common pool bound.
\end{lemma}
\begin{proof}
The mass equality in \eqref{s5:global-constants} implies
$\sum_iY_i=\sum_jV_j$, since both sides subtract the same total bypass
flow. Summing \eqref{s5:contract-quality} gives
\[
 q\sum_jV_j=\sum_j b_jB_j-\sum_{ij}C_i z_{ij}
                         =\sum_i C_iY_i.
\]
Thus actual pool mixing holds coordinate by coordinate. At positive
throughput this identifies $q$ as the feed average. At zero throughput
nonnegative incident flows all vanish, so their quality masses vanish.
Every output equation is exactly its original equation after substitution,
even when its outlet is absent. The converse follows by summing original
mass and quality conservation. A common bound has not been derived.
\end{proof}

\begin{lemma}[Signed node intervals and connected cuts]\label{s5:cuts}
For a directed graph with finite signed arc bounds $\ell_e\le u_e$
and node bounds $\alpha_v\le\beta_v$, a flow satisfying
\[
 \ell\le w\le u,\qquad
 \alpha_v\le\operatorname{div}_w(v)\le\beta_v
\]
exists if and only if, for every vertex subset $S$,
\begin{align}
 \alpha(S)&\le u(\delta^+S)-\ell(\delta^-S),\notag\\
 \beta(S)&\ge \ell(\delta^+S)-u(\delta^-S).
 \label{s5:hoffman-cuts}
\end{align}
Here divergence is outgoing minus incoming, and set arguments indicate
sums. It suffices to check connected induced subsets. A disjoint union
of paths, cycles, and isolated vertices requires $O(n^2)$ such sets,
with at most two cut edges each. Prescribing boundary arc flows gives
the same description after shifting their signed contributions out of
adjacent node intervals and retaining their arc bounds separately.
\end{lemma}
\begin{proof}
This is the signed interval transshipment criterion of
\citet[Corollary~11.2i, p.~175]{s5:schrijver2003}, with the divergence
sign reversed. To derive it from circulation feasibility, add a root
and an arc from the root to each $v$, bounded by
$[\alpha_v,\beta_v]$. Conservation at $v$ makes that arc equal its
original divergence; conservation at the root follows by summing
original divergences. The circulation inequalities on subsets excluding
the root and on complements of subsets containing it give the two
inequalities in \eqref{s5:hoffman-cuts}. Conversely those two families
are all circulation cuts, so the circulation criterion supplies the
flow. Signed bounds cause no difficulty: subtract lower bounds before
the usual nonnegative-capacity feasibility test.

If $S$ is disconnected in its induced underlying graph, both node sums
and both cut expressions add over its connected components, since no
edge joins distinct components within $S$. Their separate inequalities
therefore imply its inequalities. The connected subsets of a path are
intervals; those of a cycle are cyclic intervals and the whole cycle.
Whole components and isolated vertices must be retained. For prescribed
boundary flow $t_e$ with divergence sign $\sigma_{ve}$ at $v$, the
internal divergence lies between
$\alpha_v-\sum_e\sigma_{ve}t_e$ and
$\beta_v-\sum_e\sigma_{ve}t_e$. This is an exact substitution.
\end{proof}

\begin{theorem}[Quadratic-field feasibility]\label{s5:quadratic-field}
Suppose every source and product is exactly contracted, the quality is
scalar, and the bypass graph has maximum degree two. If the common pool
lower bound is zero and the upper bound is redundant, feasibility is
polynomial-time decidable. A feasible original flow can be constructed
in one real field of degree at most two over $\Q$, with polynomial
encoding length. The same result holds for arbitrarily many attributes
of affine input rank at most one.
\end{theorem}
\begin{proof}
Check \eqref{s5:global-constants}. If the pool has no allowed feed or no
allowed outlet, force its incident flows to zero, check their bounds,
and solve the remaining rational LP with all external contracts. Otherwise
restrict $q$ to the closed interval spanned by allowed feed qualities.
This loses no active flow by mixing; an inactive pool can be assigned
any value in that interval. At every distinct source quality in this
interval, including its endpoints, solve the original fixed-$q$ physical
LP with actual pool arcs and balances. These are rational LPs.

On each remaining open interval, every $\gamma_i=C_i-q$ is nonzero
with fixed sign. Set
\begin{equation}
 w_{ij}=\gamma_i z_{ij},\qquad z_{ij}=w_{ij}/\gamma_i.
 \label{s5:scaling}
\end{equation}
Keep the orientation input to product even when $w$ is negative. Feed
bounds $\ell_i^P\le Y_i\le u_i^P$ become a bound on
$\sum_jz_{ij}$ between $a_i-u_i^P$ and $a_i-\ell_i^P$. Multiplication
by $\gamma_i$, with order determined by its sign, gives an affine
interval on input divergence $\sum_jw_{ij}$. Every original bypass
interval likewise gives affine signed bounds on its $w_{ij}$.
Equation~\eqref{s5:contract-quality} becomes
\begin{equation}
 \sum_i w_{ij}=R_j(q),\qquad R_j(q)=b_j(B_j-q),
 \label{s5:output-divergence}
\end{equation}
so product divergence is exactly $-R_j(q)$.

The outlet interval is a bound $L_j\le s_j\le U_j$ on total bypass
$s_j$, where $L_j=b_j-u_j^P$ and $U_j=b_j-\ell_j^P$. With two inlets
of distinct qualities, write $\gamma_h=C_h-q$. Solving
$w_1+w_2=R_j$ and $s_j=w_1/\gamma_1+w_2/\gamma_2$ gives
\begin{equation}
 w_1=\frac{\gamma_1(R_j-\gamma_2s_j)}{C_1-C_2}.
 \label{s5:outlet-arc}
\end{equation}
Its coefficient of $s_j$ is
$-\gamma_1\gamma_2/(C_1-C_2)$, nonzero of fixed sign. Thus the outlet
interval is exactly the interval on $w_1$ with endpoints
\begin{equation}
 \frac{\gamma_1(R_j-\gamma_2L_j)}{C_1-C_2},\qquad
 \frac{\gamma_1(R_j-\gamma_2U_j)}{C_1-C_2}.
 \label{s5:quadratic-bounds}
\end{equation}
They are quadratic polynomials in $q$; retain
\eqref{s5:output-divergence} and all original arc bounds too.
If the qualities coincide, $\gamma s_j=R_j$ instead, so the outlet
interval gives the pure parameter condition that $R_j$ lie between
$\gamma L_j$ and $\gamma U_j$. With one inlet it gives that affine
arc interval together with $w=R_j$. With zero inlets require
$L_j\le0\le U_j$ and $R_j=0$. These cases include zero demands and
missing outlets without division by a vanishing quality difference.

Each arc now has constantly many lower and upper quadratic candidates.
Its effective bounds are the maximum lower and minimum upper candidate.
Include every lower-versus-upper consistency inequality. Node bounds
are affine. Apply Lemma~\ref{s5:cuts}. Each connected cut uses at most
two arc bounds. The first right side of \eqref{s5:hoffman-cuts} is the
minimum over all upper-outgoing/lower-incoming candidate choices; the
second is the maximum over lower-outgoing/upper-incoming choices.
Thus each cut is equivalent to the conjunction over all its constantly
many choices. The result is polynomially many univariate quadratic
inequalities and the pure parameter conditions. No capacity disjunction
or exponential subset enumeration is required.

Exact root isolation and sign tests decide their conjunction on the
open source-quality interval. Include isolated roots of these new
polynomials; only the original source-quality endpoints are excluded
from the nonsingular chart. Coefficients have polynomial bit length,
since they use constant-degree products, rational quality differences,
and polynomially many sums. A feasible open interval contains a
polynomial-bit rational sample, and a feasible isolated boundary is a
root of a quadratic polynomial with polynomial coefficient bit length. Thus choose $q$
of degree at most two.

At this $q$, evaluate all effective bounds and solve the signed
circulation in $\Q(q)$. The standard reduction in
\citet[Corollary~11.3a, p.~176]{s5:schrijver2003}, or the algebraic LP
algorithm, constructs a feasible point. Its constant incidence matrix
has vertex coordinates given by rational linear combinations of bounds
in this same field, with polynomial coefficient length. Dividing by
the nonzero $\gamma_i$ and recovering $Y_i,V_j$ preserves that field
and polynomial encoding. Lemma~\ref{s5:conservation} proves physical
feasibility. Singular LP witnesses are rational.

For affine rank one, compute $C_i=C_0+d\,\xi_i$ with rational
$d\ne0$ and rational scalars $\xi_i$. Reject a positive-demand product
outside that line, and otherwise write $B_j=C_0+d\,\bar\xi_j$.
At zero demand its quality vector may be replaced by any vector on the
line. Local mass conservation cancels the $C_0$ terms, and any nonzero
coordinate of $d$ makes all vector equations equivalent to the scalar
ones in $\xi_i,\bar\xi_j$. Rational affine algebra has polynomial
encoding and changes no physical flow bound. Rank zero is a rational
LP after checking positive-demand product compatibility.
\end{proof}

The proof also gives a reusable boundary description
\label{s5:fixed-boundary}: after retaining a fixed number of transformed
boundary flows, Lemma~\ref{s5:cuts} gives polynomially many rows linear
in those coordinates and quadratic in $q$. Capacity bounds on a removed
boundary arc remain explicit. A physical boundary coordinate $z_{ij}$
can be linked by the bilinear equation in \eqref{s5:scaling}.

\begin{example}[A common bound is an additional constraint]
\label{s5:capacity-counterexample}
Take unit supplies of qualities zero and two. The clean input bypasses
only to a product of demand one and exact quality $1/2$; the dirty input
bypasses only to a product of demand one and exact quality $3/2$.
Both inputs feed the pool and both products receive its flow; all
individual upper capacities are one and lower bounds zero. The local
contracts force
\[
 Y_0=V_{1/2}=\frac1{2q},\qquad
 Y_2=V_{3/2}=\frac1{2(2-q)},\qquad \tfrac12\le q\le\tfrac32.
\]
Hence $T=1/[q(2-q)]\ge1$. Local feasibility holds at $q=1$, but common
upper capacity $3/4$ makes the physical model infeasible. The global
identities do not imply a restrictive pool bound. The example is a
scope distinction, not a hardness result.
\end{example}

\subsection{A bounded number of contract exceptions}\label{s5:exceptions}

Designate exceptional external sets $E_I,E_J$ and put
$s=|E_I|+|E_J|$. Ordinary inputs have exact supplies; ordinary products
have exact demand and quality contracts. Exceptional inputs may have
supply intervals, and exceptional products may have demand and quality
intervals. In this subsection the bypass maximum degree is two, and
common pool bounds are redundant: their lower bound is zero and their
upper bound is at least a valid total throughput bound. For example,
the minimum of the sums of input upper supplies, feed capacities, and
outlet capacities is such a bound. Individual pool-arc intervals remain
unrestricted. Standard economics is the objective. Before the following
algorithms, a pool with no allowed feed or no allowed outlet is removed
by forcing all its incident flows to zero, checking their bounds, and
solving the remaining rational LP with every external requirement.
Thus the nonlinear constructions may assume both kinds of pool arc exist.

\begin{lemma}[A fixed-dimensional pool-quality space]
\label{s5:quality-chart}
For fixed $s,d$, suppose pool quality is restricted to a supplied rational
affine space $q=q_0+D\eta$ of dimension $d$, where the columns of $D$
are independent, and to the closed coordinate box spanned by the allowed
feed vectors. For a pool with both an allowed feed and an allowed outlet,
the preceding exception model with these quality restrictions admits polynomial-time
feasibility and standard economic optimization, with a polynomial-size
common real-algebraic original optimizer. Ambient quality dimension and
input-quality rank may grow. The retained dimension is at most $d+5s$.
\end{lemma}
\begin{proof}
Use the supplied closed coordinate box. At an active pool its balance
will imply this box; an inactive pool in this branch is required to have
a representative quality in the supplied intersection. If the
supplied affine space misses this box the branch is empty. Select $d$
independent coordinate rows of $D$. They express $\eta$ as a rational
affine function of bounded quality coordinates and hence give explicit
finite polynomial-bit bounds on $\eta$. Retain all box inequalities.

Retain each exceptional node's incident pool flow and every bypass arc
in $H$ incident to an exceptional node. There are at most $s$ such
pool flows and $2s$ bypass flows. For exceptional nodes define the core
expressions
\[
 A_i=y_i+\sum_jz_{ij},\qquad
 D_j=v_j+\sum_i z_{ij},\qquad M_j=qv_j+\sum_iC_i z_{ij}.
\]
Keep their supply, demand, quality, and arc bounds. At ordinary nodes
eliminate feeds and outlets by exact contracts as before. Preserve
\eqref{s5:contract-quality} at every ordinary product. Put in the core
one mass equation and all ambient quality equations
\begin{align}
 \sum_{i\notin E_I}a_i+\sum_{i\in E_I}A_i
 &=\sum_{j\notin E_J}b_j+\sum_{j\in E_J}D_j,\notag\\
 \sum_{i\notin E_I}C_i a_i+\sum_{i\in E_I}C_i A_i
 &=\sum_{j\notin E_J}b_jB_j+\sum_{j\in E_J}M_j.
 \label{s5:exception-balances}
\end{align}
To verify sufficiency of these equations, subtract total bypass mass
from the first equality to obtain actual pool mass balance. In each
quality coordinate, sum the ordinary output equations and the
exceptional masses $M_j$; subtracting bypass quality mass from the
second equality gives
$q\sum_jv_j=\sum_iC_i y_i$. Thus actual pool quality balance holds
for every local lift. At zero throughput all pool arcs vanish.

For each full source vector $C_i$ belonging to the supplied affine space
and box, solve the original fixed-vector-quality physical LP, including
its objective. These cover all $q=C_i$ singularities, even for inputs
without a feed. At every other $q$ some vector in the explicit list
\begin{equation}
 \beta(k)=(1,k,\ldots,k^{K-1}),\qquad
 1\le k\le |I|(K-1)+1
 \label{s5:directions}
\end{equation}
has $\gamma_i=\beta(k)\mathbin{\cdot}(C_i-q)\ne0$ for every input.
Indeed each nonzero $C_i-q$ defines a nonzero polynomial of degree at
most $K-1$ in $k$, excluding at most $K-1$ list members. The union of
all excluded sets cannot exhaust the list. Its size and coefficient
bit lengths, $O(K\log(|I|K+1))$ per entry, are polynomial even when
$K$ grows. For $K=1$ one direction suffices, and $K=0$ is an LP.

Fix a direction. Define $w_{ij}=\gamma_i z_{ij}$ and, for retained
arcs, keep both $w_{ij}$ and its bilinear link to $z_{ij}$. This gives
at most $d+s+2s+2s=d+5s$ variables. Projecting the ordinary product
equation onto $\beta$ gives
\[
 \sum_iw_{ij}=R_j,\qquad R_j=b_j\beta\mathbin{\cdot}(B_j-q).
\]
Original bypass and ordinary input bounds become affine signed arc
and divergence bounds, after partitioning by the nonzero signs of
$\gamma_i$. For a two-inlet ordinary product set
$c_h=\beta\mathbin{\cdot}C_h$. If $c_1\ne c_2$, its outlet bounds
become the quadratic arc interval of \eqref{s5:quadratic-bounds}, with
$c_1-c_2$ in the denominator and the projected $R_j$. If $c_1=c_2$,
use the corresponding pure parameter condition instead.

The projection alone does not enforce all quality coordinates. For
each ambient coordinate $h$, substitute $w_2=R_j-w_1$ into its original
quality equation and multiply by the nonzero $\gamma_1\gamma_2$.
The exact residual equation is
\begin{align}
 a_h(q)w_1&=f_h(q),\notag\\
 a_h&=(C_{1h}-q_h)\gamma_2-(C_{2h}-q_h)\gamma_1,\notag\\
 f_h&=\gamma_1\bigl[b_j(B_{jh}-q_h)\gamma_2
                         -(C_{2h}-q_h)R_j\bigr].
 \label{s5:residual-quality}
\end{align}
The quadratic terms of $a_h$ cancel, and the quadratic terms inside
the brackets of $f_h$ cancel; hence $a_h$ is affine and $f_h$ quadratic
in $\eta$. Include every $a_h$ in the affine sign partition. If
$a_h=0$, impose $f_h=0$. Otherwise impose both bounds
$w_1=f_h/a_h$ with known denominator sign. Keep all such equations,
including redundant ones, on the same chosen incident arc.

With one inlet, projected quality fixes $w=R_j$, and each additional
coordinate gives
\[
 (C_{ih}-q_h)R_j=b_j(B_{jh}-q_h)\gamma_i.
\]
Keep these parameter equalities and the outlet's affine interval on
$w$. With no inlets, require $0$ in the bypass-total interval and
$b_j(B_j-q)=0$ in every coordinate. Thus every ordinary quality
condition is preserved, including equal projected qualities and zero
coefficients. No assumption that the ambient vectors lie in the supplied
quality space was used.

The signs of the polynomially many affine functions $\gamma_i,a_h$
have polynomially many realizations in fixed dimension $d$. Keep
lower-dimensional sign conditions, and require $\gamma_i\ne0$ strictly.
On each such cell every internal arc has polynomially many rational
bound candidates, of numerator degree at most two and denominator
degree at most one. Delete exceptional nodes and remove retained
boundary arcs from the internal graph, subtracting their signed $w$
contributions from adjacent divergence intervals. A candidate bound
assigned to a boundary arc remains a core row. An arc between two
exceptional nodes requires no internal projection.

Residual components are paths, cycles, or isolated nodes. Apply
Lemma~\ref{s5:cuts} to connected sets. Each cut crosses at most two
internal edges; conjunction over their candidate combinations is still
polynomial, even with polynomial-size candidate lists. Include every
lower-versus-upper candidate comparison on each arc. Clear only the
at most two selected affine denominators in each inequality, using their
known signs. Node sums are affine in $\eta$ and linear in boundary
$w$ coordinates, so degree four is a safe bound for all resulting
polynomial inequalities. Pure parameter conditions and retained
candidate rows have the same bound. Multiplying all internal-arc
denominators would be unnecessary and would lose this constant-degree
property.

The resulting formula has fixed dimension, constant degree,
polynomially many rows, and polynomial coefficient length. Together
with \eqref{s5:exception-balances} and the retained physical rows it is
exactly the original model on the chart. Standard profit reduces to
\begin{equation}
 \sum_{j\notin E_J}R_j^{\rm econ}b_j-
 \sum_{i\notin E_I}c_i a_i+
 \sum_{j\in E_J}R_j^{\rm econ}D_j-
 \sum_{i\in E_I}c_i A_i,
 \label{s5:exception-profit}
\end{equation}
where $R_j^{\rm econ}$ denotes revenue and is unrelated to the
projected quality residual $R_j$. It is affine in original retained
flow coordinates.

Use quantifier elimination to describe the attainable profit values
on every actual sign cell, and unite these descriptions with the
fixed-source-vector LP value intervals. The union is precisely the
attainable set of the original model with the closed affine quality
restriction and box, a compact set. If nonempty it has an attained
maximum. Select that value and sample a point on a cell attaining it.
Do not close a chart at a vanishing denominator: its polynomial rows
can lose information there, while the singular LPs retain it.

At the sample, all residual incidence bounds belong to its common
represented real-algebraic field. Solve each residual signed flow LP
in that field, then divide by nonzero $\gamma_i$ and recover ordinary
feeds and outlets from their contracts. A constant incidence matrix
has vertex coordinates that are rational linear combinations of its
bounds, and the sample has polynomial degree and encoding. Therefore
recovery remains polynomial in the same field. Equations
\eqref{s5:exception-balances} prove the actual pool balances for every
recovered flow.
\end{proof}

\begin{corollary}[Scalar and fixed-rank exceptions]\label{s5:fixed-rank}
For fixed $s$ and fixed affine input-quality rank $t$, the exception
model admits polynomial-time feasibility and standard economic
optimization. For scalar quality the core dimension is at most $1+5s$
and every chart row may be chosen quadratic. For fixed rank $t$ the
bound is $t+5s$. Neither assertion gives a quadratic-field witness
bound when exceptions are present.
\end{corollary}
\begin{proof}
For scalar quality the above chart argument uses only the signs of
$C_i-q$, the scalar intervals \eqref{s5:quadratic-bounds}, the boundary
cut description, and the two equations
\eqref{s5:exception-balances}; there are no residual vector equations
or variable denominators in arc candidates. Thus all rows are quadratic
and dimension is $1+5s$.

For affine rank $t$, compute rational coordinates
$C_i=C_0+D\xi_i$ with $t$ independent columns. A positive-demand
ordinary product must belong to this affine hull; reject it otherwise.
Zero-demand qualities are immaterial. Express every original exceptional
quality inequality in the coordinate quality masses and throughput;
keep all such inequalities. Local mass conservation cancels $C_0$, so
coordinate quality equations are equivalent to all ambient equations.
Apply Lemma~\ref{s5:quality-chart} in $t$ coordinates, or take its
supplied space to be this full affine hull. Rank zero is a fixed-quality
LP. All coordinate and bound computations are polynomial-bit rational
linear algebra.
\end{proof}

\begin{theorem}[Arbitrary-quality bounded exceptions]
\label{s5:arbitrary-exceptions}
For fixed $s$, the exception model admits polynomial-time exact
feasibility and standard economic optimization with a polynomial-size
common real-algebraic original optimizer. The number of quality
coordinates and their affine input rank are unrestricted.
\end{theorem}
\begin{proof}
Consider an ordinary product $j$ with positive pool inflow $v_j>0$.
Its nonnegative flow and exact contract imply $b_j>0$, and
\[
 q=\frac{b_j}{v_j}B_j-\sum_i\frac{z_{ij}}{v_j}C_i.
\]
The coefficients sum to one by product mass conservation. Since at most
two bypass inlets exist,
\begin{equation}
 q\in\operatorname{aff}\bigl(\{B_j\}\cup
                    \{C_i:ij\text{ is a bypass}\}\bigr),
 \label{s5:active-space}
\end{equation}
an affine space of dimension at most two. Enumerate this space for each
positive-demand ordinary product and apply
Lemma~\ref{s5:quality-chart} with $d\le2$. Rational Gaussian elimination
gives polynomial-bit bases and bounds. Dimension zero is a physical
fixed-quality LP. There are polynomially many spaces. No branch needs
to require that its designated outlet actually be positive: every
branch retains all original constraints and is sound; every solution
with an active ordinary outlet belongs to a branch by
\eqref{s5:active-space}.

It remains to handle flows whose only active outlets are exceptional.
Force all ordinary outlets to zero, checking their lower bounds.
Let $J_E\subseteq E_J$ be the products with an allowed pool outlet and
put $r=|J_E|$. For $j\in E_J\setminus J_E$, put $v_j=\theta_j=0$
by convention. If $r=0$, force all pool arcs to zero and solve the
remaining rational LP. Otherwise retain exceptional input feeds, the
$r$ outlet flows indexed by $J_E$, every bypass arc incident to an
exceptional node, and $r$ outlet fractions $\theta_j$, $j\in J_E$.
There are at most $3|E_I|+4|E_J|$ coordinates. Exceptional source
withdrawals $A_i=y_i+\sum_jz_{ij}$ are core expressions. Ordinary
inputs give the local affine feed bounds, while ordinary products now
give bypass-only exact mass and vector-quality equations. Delete the
exceptional nodes and apply Theorem~\ref{s4:path-projection} to all
remaining rational path relations; solve detached components by LP.

Define core expressions
\begin{align}
 S&=\sum_{i\notin E_I}a_i+\sum_{i\in E_I}A_i,&
 G&=\sum_{i\notin E_I}C_i a_i+\sum_{i\in E_I}C_i A_i,\notag\\
 T&=S-\sum_{j\notin E_J}b_j-
                     \sum_{ij:\,j\in E_J}z_{ij},&
 Q&=G-\sum_{j\notin E_J}b_jB_j-
                     \sum_{ij:\,j\in E_J}C_i z_{ij}.
 \label{s5:exceptional-only-masses}
\end{align}
Every term uses only retained coordinates or constants. Summing
ordinary product contracts and all source equations proves
$T=\sum_i y_i$ and $Q=\sum_iC_i y_i$ for every local lift. Impose
$\theta_j\ge0$ and $v_j=\theta_jT$ for $j\in J_E$, and
$\sum_{j\in J_E}\theta_j=1$. Keep the demand and homogeneous quality
bounds at every $j\in E_J$, including products without pool outlets,
using quality mass
$\theta_jQ+\sum_iC_i z_{ij}$ and throughput $v_j+\sum_i z_{ij}$,
as well as all retained arc and source bounds. These rows are quadratic
in a fixed number of variables independently of the ambient quality
count. At positive $T$ they give actual pool quality $Q/T$. At zero
$T$, all nonnegative feeds vanish and $Q=0$, so fraction choices are
harmless. Conversely every physical flow of this branch supplies its
fractions, or any simplex point if inactive. Thus this branch is exact,
closed, bounded, and polynomially described.

Standard profit is again \eqref{s5:exception-profit}, hence depends
only on the core. Optimize this branch and all active-space branches
and select the largest feasible value. They cover every physical flow;
each feasible branch is compact before its internal chart partition,
so an optimum is attained. The quality branches use incidence recovery;
the exceptional-only branch uses rational endpoint lifting and local
feed recovery. Both yield one common field of polynomial degree and
encoding. This proves the theorem.
\end{proof}

A fixed number of designated additional arc costs can be included by
declaring their external endpoints exceptional and preserving their
original contracts there. Their actual arc flows are then retained.
The increase in $s$ is fixed, so the proof applies with these affine
core costs. This is not unrestricted dense arc-cost optimization.
The common-bound assumption in this subsection also remains necessary
for the argument: total pool throughput need not be a core coordinate.
Theorem~\ref{s3:five-thm} gives the complementary strong feasibility
hardness already with five exceptions when ordinary products may have
three bypass inlets. In particular the positive bypass-degree-two
condition here is consistent with the total-output-degree-three
hardness boundary and with the positive contracts used there.

\subsection{Restrictive common pool bounds}\label{s5:common-capacity}

The signed-flow feasibility cuts do not constrain the total unscaled
bypass flow. To recover that quantity, we compute its extreme values
as a linear objective on node divergences. Two supporting facts make
that computation polynomial over the quality parameter. The classical
box-truncated submodular rank formula supplies a support value at a
fixed quality; the additional step is to compute all prefix rank values
uniformly as the quality changes. The binary path recurrence below
provides that step without enumerating all bypass flow bases.

\begin{lemma}[Symbolic binary path and cycle minimization]
\label{s5:clamp}
Let $\eta$ have fixed dimension, and consider binary path labels
$x_1,\ldots,x_n$ with energy
\[
 \sum_iU_i(x_i;\eta)+\sum_{i=2}^n
 \bigl[a_i(\eta)\mathbf1_{(x_{i-1},x_i)=(0,1)}+
       b_i(\eta)\mathbf1_{(x_{i-1},x_i)=(1,0)}\bigr].
\]
Suppose all terms are rational polynomials of fixed degree and
$a_i,b_i\ge0$ on the specified domain; unary terms may be signed.
The minimum has an explicit polynomial formula on each of polynomially
many realizable polynomial sign conditions, with polynomial coefficient
length and the same fixed degree. The same conclusion holds for cycles.
\end{lemma}
\begin{proof}
Let $E_i(t)$ be the minimum prefix energy ending in label $t$, put
$D_i=E_i(1)-E_i(0)$, and $u_i=U_i(1)-U_i(0)$. The two ordinary
prefix recurrences give
\begin{align*}
 E_i(0)&=E_{i-1}(0)+U_i(0)+\min(0,D_{i-1}+b_i),\\
 D_i&=u_i+\min(D_{i-1},a_i)-\min(0,D_{i-1}+b_i)\\
    &=u_i+\operatorname{clamp}(D_{i-1},-b_i,a_i),
 \qquad D_1=u_1.
\end{align*}
The last identity follows by the three cases $D<-b_i$,
$-b_i\le D\le a_i$, and $D>a_i$; nonnegative $a_i,b_i$ put the
endpoints in the stated order. Set $P_i=\sum_{h=1}^iu_h$ and
$Z_i=D_i-P_i$. Then
\begin{equation}
 Z_1=0,\qquad Z_i=\operatorname{clamp}
          (Z_{i-1},-b_i-P_{i-1},a_i-P_{i-1}).
 \label{s5:clamp-recurrence}
\end{equation}
Every $Z_i$ is selected from zero and the $2(i-1)$ preceding clamp
endpoints. Form the common list of those polynomials and all $-P_i$.
Signs of its pairwise differences determine every clamp selection,
every increment $\min(0,D_{i-1}+b_i)$, and the final comparison between
$E_n(0)$ and $E_n(1)$. On each realizable sign condition, execute the
recurrence symbolically; its value is a sum of polynomially many input
terms of the original degree. Equal values can be resolved by either
choice, including on lower-dimensional conditions. Coefficient bit lengths
remain polynomial. A polynomial-size list of fixed-degree polynomials
has polynomially many realizable sign conditions in fixed dimension,
constructible by the cited real-algebraic algorithms.

For a cycle, fix $x_1$ to zero and then to one. The edge from the last
vertex to the fixed first vertex is a unary cost at the last vertex;
the first edge is handled by the fixed initial condition (or absorbed
at the second vertex). Apply the path construction and compare the two
value formulas on their common sign refinement. A product of two
polynomial bounds is polynomial. This proof is specific to the binary
submodular energy and its clamp recurrence.
\end{proof}

\begin{lemma}[Box-bounded divergence support]\label{s5:box-base}
Let $\ell\le u$ be finite signed arc bounds on a directed graph with
vertex set $V$, and let $\alpha\le\beta$ be finite node bounds.
Assume
\[
 \mathcal D=\{\operatorname{div}(w):\ell\le w\le u,
                  \ \alpha\le\operatorname{div}(w)\le\beta\}
\]
is nonempty. Define
\begin{equation}
 f(S)=u(\delta^+S)-\ell(\delta^-S),\qquad
 g(S)=\min_{T\subseteq V}\{f(T)+\beta(S\setminus T)
                                     -\alpha(T\setminus S)\}.
 \label{s5:box-rank}
\end{equation}
Then $g$ is submodular, $g(\varnothing)=g(V)=0$, and
\[
 \mathcal D=B(g):=\{d:d(S)\le g(S)\ (S\subseteq V),\ d(V)=0\}.
\]
If $c_1\ge\cdots\ge c_n$ is any ordering of a real cost vector and
$S_k=\{1,\ldots,k\}$, then
\begin{equation}
 \max_{d\in\mathcal D}c^Td
   =\sum_{k=1}^{n-1}(c_k-c_{k+1})g(S_k),\qquad
 d_k^*=g(S_k)-g(S_{k-1})
 \label{s5:greedy-support}
\end{equation}
gives an attaining divergence vector. Minimum support is obtained by
negating the maximum for $-c$.
\end{lemma}

The box rank formula and greedy rule are classical. They are stated in
\citet[Theorems~1--2, equation~(11), p.~192]{s5:shioura2013} for the
maximal vectors of a box-truncated submodular polyhedron. We prove the
needed base-polyhedron specialization, including its total-sum condition.

\begin{proof}
Before the node box is imposed, feasible divergences are exactly $B(f)$.
Necessity follows by summing divergences. Conversely, prescribe a vector
$d$ with $d(V)=0$ and $d(S)\le f(S)$; its complementary inequality is
$d(S)\ge\ell(\delta^+S)-u(\delta^-S)$. Lemma~\ref{s5:cuts} with
$\alpha=\beta=d$ supplies the flow. Also
\[
 f(T)=\sum_{e\in\delta^+T}(u_e-\ell_e)
                         +\sum_{v\in T}\operatorname{div}(\ell)_v.
\]
A nonnegative directed cut function is submodular, as verified edge
by edge on the four membership possibilities, and the second term is
modular. Thus $f$ is submodular even when $\ell,u$ are signed, and
$f(\varnothing)=f(V)=0$.

The expression minimized in \eqref{s5:box-rank} is jointly submodular
in $(S,T)$. Its node term takes values $0,\beta_v,-\alpha_v,0$ at
membership pairs $(0,0),(1,0),(0,1),(1,1)$; the only nontrivial
submodularity inequality is $\alpha_v\le\beta_v$. Add the submodular
term $f(T)$. For minimizing pairs $(S_1,T_1),(S_2,T_2)$, the product
lattice inequality bounds their sum below by the values at union and
intersection, which in turn bound $g(S_1\cup S_2)+g(S_1\cap S_2)$.
Hence partial minimization gives submodular $g$.

For every $d\in\mathcal D$ and every $S,T$,
\[
 d(S)=d(T)+d(S\setminus T)-d(T\setminus S)
       \le f(T)+\beta(S\setminus T)-\alpha(T\setminus S).
\]
Thus $d(S)\le g(S)$. At $S=\varnothing,V$, this inequality and the
choices $T=\varnothing,V$ respectively prove
$g(\varnothing)=g(V)=0$, using nonemptiness. This gives
$\mathcal D\subseteq B(g)$. Conversely, choose $T=S$ to see
$g(S)\le f(S)$, so $B(g)\subseteq B(f)$. At $S=\{v\}$ choose
$T=\varnothing$ to get $d_v\le\beta_v$; at $S=V\setminus\{v\}$
choose $T=V$ to get $d(V\setminus\{v\})\le-\alpha_v$, hence
$d_v\ge\alpha_v$. This proves equality of the two polyhedra.

For completeness, the greedy vector is feasible by diminishing
increments. For any subset $A$, and each $k\in A$, the marginal of
adding $k$ to $S_{k-1}$ is at most its marginal when added to
$A\cap S_{k-1}$, since the latter set is contained in the former.
Summing these inequalities gives $\sum_{k\in A}d_k^*\le g(A)$.
The full sum telescopes to $g(V)=0$. For any feasible $d$,
\[
 c^Td=c_n d(V)+\sum_{k=1}^{n-1}(c_k-c_{k+1})d(S_k)
       \le\sum_{k=1}^{n-1}(c_k-c_{k+1})g(S_k).
\]
The greedy vector attains equality at each prefix. This handles signed
costs and ties, and proves \eqref{s5:greedy-support}.
\end{proof}

\begin{theorem}[Common lower and upper throughput bounds]
\label{s5:restrictive-capacity}
The fully contracted scalar or affine-rank-one model of
Theorem~\ref{s5:quadratic-field} admits polynomial-time feasibility with
an arbitrary common throughput interval $[L_P,U_P]$. All individual
arc bounds, including positive lower bounds, remain. A feasible physical
flow is constructible in one real algebraic field of polynomial degree
and encoding length. No quadratic-field bound is asserted here.
\end{theorem}
\begin{proof}
Intersect the common interval with nonnegativity and a valid finite
throughput bound; reject an empty intersection. If the pool has no feed
or no outlet, force all incident flows to zero and explicitly require
$L_P=0$, check incident bounds, and solve the remaining original LP.
Check \eqref{s5:global-constants}; treat every singular source quality
by its original rational physical LP, now retaining $[L_P,U_P]$.

On a nonsingular quality interval perform the signed transformation
of Theorem~\ref{s5:quadratic-field}. Split at roots of all arc-candidate
comparisons. There are polynomially many quadratic comparisons and
hence polynomially many interval and singleton cells. On each cell the
effective $\ell_e,u_e$ are explicit quadratic polynomials; node
$\alpha,\beta$ are affine. Keep every pure parameter condition and
connected cut. These enforce $\mathcal D(q)\ne\varnothing$ before
Lemma~\ref{s5:box-base} is applied.

Let $A=\sum_i a_i$ and assign costs
\[
 c_i(q)=\frac1{C_i-q}\quad\hbox{at inputs},\qquad
 c_j(q)=0\quad\hbox{at products}.
\]
Since input divergence is the sum of its outgoing transformed bypasses,
\begin{equation}
 Z:=\sum_{ij}z_{ij}=c(q)^T\operatorname{div}(w),\qquad T=A-Z.
 \label{s5:throughput-support}
\end{equation}
The cost order is fixed on each nonsingular source-quality interval:
input cost differences have numerator $C_h-C_i$ and denominator
$(C_i-q)(C_h-q)$, both of known sign; the comparison with zero is
also fixed. Repeated input qualities merely give ties.

For every prefix in the orders for $c$ and $-c$, evaluate $g(S)$
symbolically. Its definition \eqref{s5:box-rank}, with membership of
$T$ as binary labels, is a path/cycle energy. Indeed $f(T)$ is a directed
cut with capacities $u_e-\ell_e\ge0$ plus the unary terms
$\operatorname{div}(\ell)_v\mathbf1_{v\in T}$; the remaining box
terms are unary too. All these costs are quadratic polynomials on the
cell. Lemma~\ref{s5:clamp} gives an explicit quadratic minimum on
polynomially many univariate sign cells. Different connected components
minimize independently, including isolated vertices. Take the union
of all univariate breakpoints over all prefixes, orders, and components.
This remains polynomial; it is not a product enumeration of their
partitions. Keep singleton and equality cells.

On the resulting cells, \eqref{s5:greedy-support} supplies rational
functions $Z_{\min}(q),Z_{\max}(q)$ with polynomial degrees and
coefficient lengths. A common denominator is the product of the
distinct nonzero factors $C_i-q$; its degree is at most the input count
and its sign is fixed. Combining polynomially many quadratic numerators
with these factors increases degree and coefficient bit length only polynomially.
For a fixed feasible $q$, compactness and convexity of the signed flow
polytope imply that its attainable bypass totals form exactly the closed
interval $[Z_{\min},Z_{\max}]$. It meets $[A-U_P,A-L_P]$ exactly when
\begin{equation}
 Z_{\min}(q)\le A-L_P,\qquad Z_{\max}(q)\ge A-U_P.
 \label{s5:capacity-test}
\end{equation}
Clear the known-sign denominator. These are univariate polynomial
inequalities of polynomial degree and bit length. Root isolation and
sign testing on each actual cell, together with retained cut and local
conditions, decide feasibility in polynomial bit time.

Choose a feasible rational interval sample or an algebraic boundary
point $q$ of polynomial degree and encoding. For each support extremum,
Lemma~\ref{s5:box-base} constructs an attaining divergence by differences
of the computed prefix rank values. These values belong to $\Q(q)$
with polynomial encoding. Realize each divergence by a bounded-flow LP
in the same field. Its incidence matrix has a polynomial-encoding
vertex; the cited algebraic LP algorithm constructs one. This also
follows from the usual lower-bound shift and a polynomial-augmentation
flow algorithm using exact comparisons in the represented field.

Choose $Z^*$ in the nonempty intersection with $[A-U_P,A-L_P]$, for
example $\max(Z_{\min},A-U_P)$. If the extrema differ, interpolate the
two realizing flows with coefficient
$(Z^*-Z_{\min})/(Z_{\max}-Z_{\min})$. If equal, use either. This
coefficient is one quotient in $\Q(q)$ of polynomial encoding.
Convexity preserves all signed arc and node bounds and makes the total
exactly $Z^*$. Divide by $C_i-q$, recover physical pool arcs, and use
Lemma~\ref{s5:conservation} to verify all balances. The resulting
throughput is in the common interval. The rank-one compression of
Theorem~\ref{s5:quadratic-field} changes no flow or capacity and therefore
extends the conclusion to affine input rank one.
\end{proof}

\begin{corollary}[Exact pool-throughput optimization]
\label{s5:throughput-optimization}
In Theorem~\ref{s5:restrictive-capacity}, both the minimum and maximum
pool throughput, and therefore standard economic cost plus a rational
linear pool-throughput cost, can be optimized exactly in polynomial bit
time with a polynomial-size common-field physical optimizer.
\end{corollary}
\begin{proof}
On each nonsingular cell retain a second variable $T$ and impose
\[
 L_P\le T\le U_P,\qquad
 A-Z_{\max}(q)\le T\le A-Z_{\min}(q).
\]
Clear the same known-sign denominators. These rows and the retained
local feasibility conditions give a formula of polynomial size and
degree in two variables. Every satisfying $(q,T)$ is physically
realizable by the extrema interpolation in the preceding proof.
Conversely every physical flow in the cell satisfies them. At each
singular quality, the original rational LP gives the full attainable
throughput interval by minimizing and maximizing $T$.

The union of all these actual-cell descriptions and singular intervals
is exactly the attainable $(q,T)$ set of the original compact physical
model with the closed feed-quality interval. Quantifier elimination
projects this polynomial-size union onto $T$ and supplies its attained
endpoints. Sample an attaining $(q,T)$ in its actual cell and use the
same interpolation to recover an original optimizer, or use the
singular LP. Fixed dimension and polynomial degrees give polynomial bit
time and a polynomial-degree common field. Exact source and product
throughputs make standard economics constant, so adding any rational
multiple of $T$ selects one of these extrema (or any feasible flow when
its coefficient is zero). No other dense arc-cost objective is included.
\end{proof}

\subsection{Two full input-quality vectors}\label{s5:two-vectors}

A different reduction permits arbitrary bypass topology and arbitrary
source supply intervals. It uses the fact that all nonlinear outlet
upper bounds have the same concave quadratic term. For interior pool quality, a single
auxiliary variable reduces the nonlinear feasibility test to one rational
convex quadratic program. Its rational optimizer yields rational
physical flows, a stronger witness conclusion than the algebraic
bounds for the degree-two algorithms.

\begin{theorem}[Two-vector feasibility with variable supplies]
\label{s5:two-vector-theorem}
Suppose the input qualities take at most two distinct full rational
vectors. Products have exact demands and exact quality vectors, sources
may have arbitrary finite rational supply intervals, and the bypass
graph is unrestricted. Individual feed and bypass lower bounds may be
positive. Every pool-outlet lower bound and the common pool lower bound
must be zero; their finite upper bounds may be restrictive. Feasibility
is polynomial-time decidable and, when feasible, has a constructible
rational original-flow witness of polynomial encoding length. This
statement concerns feasibility, including variable procurement, without
an economic or arbitrary arc-cost optimization claim.
\end{theorem}
\begin{proof}
If there are no inputs, all flows are zero and the original bounds
and homogeneous contracts can be checked directly. With one distinct
input vector, every active mixture has that vector,
so checking compatible positive-demand product qualities leaves a
rational LP. With two distinct vectors, choose a coordinate on their
rational affine line that assigns the vectors the scalar values zero
and one. Every positive-demand product vector must lie on the segment
between them; otherwise it cannot be a convex mixture and the instance
is infeasible. Its scalar coordinate is $B_j\in[0,1]$. A zero-demand
product has all incoming flows zero, so its stated quality is immaterial
and may be assigned any coordinate on the line. Local mass equations
cancel the affine offset, making scalar quality equivalent to every
original vector equation. All this is polynomial-bit rational linear
algebra. Write $I_0,I_1$ for the two input classes.

If the pool lacks a feed or an outlet, force all its arcs to zero,
check their bounds, and solve the remaining original LP with source
intervals and product contracts. Otherwise the only active scalar
qualities are in $[0,1]$. Solve the original physical LPs at $q=0$ and
$q=1$, retaining every pool balance and capacity. An inactive pool can
also be represented at either endpoint. Only the nonsingular case
$0<q<1$ remains.

Let actual withdrawal at input $i$ be $A_i$ and feed flow $y_i$.
Exact product contracts fix the total consumption of each class:
\begin{equation}
 A_0^{\rm tot}=\sum_jb_j(1-B_j),\qquad
 A_1^{\rm tot}=\sum_jb_jB_j.
 \label{s5:class-totals}
\end{equation}
These identities follow by total mass and scalar-quality conservation;
individual withdrawals need not be fixed. Put $\gamma_i=-q$ on $I_0$
and $\gamma_i=1-q$ on $I_1$, and introduce
\begin{equation}
 w_{ij}=\gamma_i z_{ij},\qquad t_i=\gamma_i A_i,
 \qquad h_i=\gamma_i y_i.
 \label{s5:two-scaling}
\end{equation}
A source interval $[L_i,U_i]$ becomes the interval with endpoints
$\gamma_iL_i,\gamma_iU_i$ on $t_i$, in the sign order fixed by its
class. Feed intervals give the corresponding bounds on $h_i$, with
$h_i=0$ for an absent feed. Bypass intervals give affine bounds on
$w_{ij}$. This retains positive feed and bypass lower bounds. Impose
\begin{equation}
 t_i=h_i+\sum_jw_{ij},\qquad
 \sum_{i\in I_0}t_i=-q A_0^{\rm tot},\qquad
 \sum_{i\in I_1}t_i=(1-q)A_1^{\rm tot}.
 \label{s5:two-input-rows}
\end{equation}
Every row is affine in the transformed variables and $q$.

For each product define
$W_{0j}=\sum_{i\in I_0}w_{ij}$ and
$W_{1j}=\sum_{i\in I_1}w_{ij}$. Its exact quality equation after
eliminating $v_j=b_j-\sum_i z_{ij}$ is
\begin{equation}
 W_{0j}+W_{1j}=b_j(B_j-q).
 \label{s5:two-output}
\end{equation}
Use $\sum_i z_{ij}=-W_{0j}/q+W_{1j}/(1-q)$ in this identity to get
\begin{equation}
 W_{0j}=q b_j(B_j-1)+q(1-q)v_j.
 \label{s5:two-outlet-identity}
\end{equation}
This uses no bound on the number of bypass inlets and remains valid
when an input class is absent at a product. Thus the outlet bounds
$0\le v_j\le U_j^P$ are exactly
\begin{equation}
 q b_j(B_j-1)\le W_{0j}\le
        q b_j(B_j-1)+U_j^P q(1-q).
 \label{s5:concave-outlet}
\end{equation}
An absent outlet is $U_j^P=0$. Summing
\eqref{s5:two-outlet-identity} gives the common upper bound $T\le U_P$
as
\begin{equation}
 \sum_jW_{0j}\le q\sum_jb_j(B_j-1)+U_Pq(1-q).
 \label{s5:concave-common}
\end{equation}
All upper coefficients $U_j^P,U_P$ are nonnegative.

We check the reverse physical correspondence before solving these
inequalities. At interior $q$, divide \eqref{s5:two-scaling} by its
nonzero class factor to recover $A_i,y_i,z_{ij}$. Their original bounds
and input conservation follow from the affine interval and local rows.
The class equations \eqref{s5:two-input-rows} give
$\sum_i A_i=A_0^{\rm tot}+A_1^{\rm tot}=\sum_jb_j$, hence recovered
pool intakes and outlets have the same total. Also
\[
 \sum_i t_i=-q A_0^{\rm tot}+(1-q)A_1^{\rm tot}
       =\sum_jb_j(B_j-q)=\sum_{ij}w_{ij}.
\]
Summing the input equations therefore gives $\sum_i h_i=0$, namely
\[
 -q\sum_{i\in I_0}y_i+(1-q)\sum_{i\in I_1}y_i=0.
\]
This is precisely the pool quality balance. The product equations give
actual product quality, and \eqref{s5:concave-outlet} gives actual
nonnegative outlet flow and its upper bound by
\eqref{s5:two-outlet-identity}. Equation~\eqref{s5:concave-common}
gives the shared capacity. At zero throughput every nonnegative pool
flow vanishes. Conversely every original interior-quality solution
satisfies all transformed rows, by these same identities.

Introduce one auxiliary variable $r$ and replace every $q(1-q)$ on
the right sides of \eqref{s5:concave-outlet} and
\eqref{s5:concave-common} by $q-r$. Let $P$ be the resulting rational
polytope of all affine rows together with $0\le q,r\le1$. It is
bounded: source, feed, and bypass capacities bound $t,h,w$ uniformly
since $|\gamma_i|\le1$. Empty $P$ is rejected. The original transformed
system at an interior $q$ is feasible if and only if some point in $P$
satisfies
\begin{equation}
 F(q,r,t,h,w):=q^2-r\le0.
 \label{s5:single-qp}
\end{equation}
Forward take $r=q^2$. In reverse, $r\ge q^2$ makes every replaced
upper bound at least as restrictive as the original one, since every
multiplying capacity is nonnegative. The lower bounds are affine and
were unchanged. Thus the equivalence holds.

The rational objective $F$ has positive semidefinite Hessian of rank
one. Exact rational convex quadratic programming over a rational
polytope is polynomial-time solvable, including an exact optimal value
and a rational optimizer of polynomial encoding length. This is the
classical algorithm of \citet[pp.~1320--1323; numbered paragraphs~3
and~6 for rational optimal solutions and recovery]{s5:kozlov1980}.
It is a convex-QP invocation, not an appeal to generic exact conic
feasibility. For the rationality assertion specifically, at a minimizer
choose its active affine rows. The corresponding KKT stationarity,
active equalities, primal inequalities, and nonnegative multipliers
form a rational linear system; every solution is optimal by convexity.
A rational basic feasible solution of that system has polynomial
encoding by determinant bounds. The algorithm constructs such an
optimizer without enumerating active sets.

The closed polytope may contain spurious transformed endpoint points,
so interior quality must still be enforced exactly. First maximize
$\epsilon$ over $P$ with
$0\le\epsilon\le q$ and $\epsilon\le1-q$. If infeasible or its optimum
is zero, $P$ contains no interior-quality point. Otherwise retain a
polynomial-bit rational point $p\in P$ with $0<q_p<1$. Now minimize
$F$ over $P$, obtaining a rational optimizer $x^*$ and value $m$.
If $m>0$, there is no solution. If $m=0$, all minimizers have the same
$q$: for two minimizers $x,y$,
\[
 F((x+y)/2)=\tfrac12F(x)+\tfrac12F(y)-\tfrac14(q_x-q_y)^2,
\]
so different qualities would contradict optimality. Thus this branch
is feasible precisely when $0<q_{x^*}<1$; then $x^*$ is a witness.
If $m=-\rho<0$, use $F(p)\le1$ (from $0\le q,r\le1$) and choose
\[
 \lambda=\frac{\rho}{2(\rho+1)},\qquad
 x=(1-\lambda)x^*+\lambda p.
\]
This rational point has interior quality, and convexity gives
$F(x)\le-\rho+\lambda(\rho+1)=-\rho/2<0$. All numbers have polynomial
encoding. This handles exact zero minima and isolated feasible interior
qualities without a numerical margin or a closure assumption.

Finally divide by the nonzero rational class factors to recover
$A_i,y_i,z_{ij}$ and set $v_j=b_j-\sum_i z_{ij}$. Every recovered flow
is rational with polynomial bit length; the established reverse mapping
verifies all physical rows. Endpoint branches already use the original
rational LPs. Returning to the original affine quality line restores
all full-vector quality equations and proves the theorem.
\end{proof}

\begin{corollary}[Exact supplies and additional outlet resources]
\label{s5:two-vector-extensions}
Theorem~\ref{s5:two-vector-theorem} includes exact supplies as a special
case. For exact supplies standard input costs and output revenues are
constant. In either supply model it also permits any polynomial number
of additional outlet rows $\sum_j\alpha_jv_j\le U$ with rational
$\alpha_j$ and $U\ge0$.
\end{corollary}
\begin{proof}
Exact supplies fix $A_i=a_i$, so the class totals require
$\sum_i a_i=\sum_jb_j$ and $\sum_{i\in I_1}a_i=\sum_jb_jB_j$.
The variables $t_i$ in \eqref{s5:two-input-rows} may then be substituted
by $\gamma_i a_i$. Eliminating $h_i$ gives the affine signed bypass
interval at each source, preserving the exact-source construction and
all its lower bounds. Economic constancy follows from fixed source and
product throughputs.

For an additional resource row, multiply
\eqref{s5:two-outlet-identity} by $\alpha_j$ and sum. Since $q(1-q)>0$,
the row is exactly
\[
 \sum_j\alpha_jW_{0j}\le
 q\sum_j\alpha_jb_j(B_j-1)+Uq(1-q).
\]
Replace its last term by $U(q-r)$ in $P$. Nonnegative $U$ gives the
same two-way equivalence with \eqref{s5:single-qp}. The signs of the
$\alpha_j$ do not change the identity or this comparison. Endpoint LPs
retain the original resource rows. The proof and rational recovery are
otherwise unchanged.
\end{proof}

Positive outlet lower bounds would give
$W_{0j}\ge qb_j(B_j-1)+\ell_j^P q(1-q)$, whose quadratic term has the
opposite inequality direction. A positive common pool lower bound has
the same obstruction. Neither is covered by this convex reduction;
they are permitted by the different, degree-two algorithm in
Theorem~\ref{s5:restrictive-capacity}. Variable source withdrawals also
make input procurement costs nonconstant, so feasibility in
Theorem~\ref{s5:two-vector-theorem} does not establish their optimization.
The restriction is two distinct complete quality vectors, not two
quality coordinates or a fixed number of physical inputs.

The established ingredients in this section are planar elimination,
difference constraints, signed circulation, submodular base optimization,
convex quadratic programming, and fixed-dimensional real algebra. Their
specializations above retain the physical contract assumptions. The
fixed-input no-bypass algorithm of \citet[Theorem~1]{s4:boland2017}
and the fixed-demand model without feed availability and common pool
capacity constraints in \citet[Assumption~2.2]{s3:baltean2018} have
different hypotheses. They do not supply the unbounded-input,
contracted-bypass classifications proved here. The additional results
are the physical transformations and their uniform parameter analysis:
conserved boundary masses, quadratic signed arc bounds, low-dimensional
quality charts for exceptions, symbolic throughput support, and the
shared-quadratic reduction for two full vectors. Each proof identifies
which contracts and cost terms its transformation preserves. This
comparison is scoped to the cited statements and is not a claim of
exhaustive priority.
